%! TEX root = ../main.tex
\chapter{Summary}

In many scientific problems, the quantity one wants to estimate is not directly observable. Instead, the data is a noisy, indirect image of it, produced by a known forward transformation. Recovering the unknown from the observed data requires inverting that transformation: a statistical inverse problem. A canonical example is Wicksell's stereological problem: spherical particles of unknown size are embedded in an opaque material, and the only observable data are the radii of the circular cross-sections left by a cutting plane, which are a known geometric transform of the true sphere radii. The same mathematical structure appears in astronomy, deconvolution problems in epidemiology, and, in a different guise, in causal inference with unmeasured confounders.

Two difficulties make these problems hard. The first is \textit{ill-posedness}: the forward operator may smooth the data, so inverting it amplifies estimation errors; naive unconstrained estimators are erratic and often have infinite variance. Exploiting \textit{shape constraints} like monotonicity, in some of the problems treated here, is the natural regularisation. The second difficulty is \textit{uncertainty quantification} (UQ). Even after constructing a well-behaved estimator, building valid confidence sets is non-trivial: the limiting distributions depend on unknown local properties of the distribution $F_0$ of the unknown (the density of the observations $g_0(x)$ and the local H\"{o}lder smoothness $\gamma$ of $F_0$), which are themselves hard to estimate. This makes naive plug-in confidence intervals unreliable.

The main approach developed in this thesis is to construct Bayesian nonparametric posteriors and then isotonise the posterior draws --- projecting each draw onto the set of monotone functions via the least concave majorant. The result, the \textit{Isotonized Inverse Posterior} (IIP), is computationally inexpensive and, as we show, automatically valid for UQ without requiring the estimation of any nuisance parameter.

Chapter~2 studies Wicksell's problem from a frequentist perspective. We prove that the \textit{isotonic inverse estimator} (IIE) --- the right-derivative of the least concave majorant of a naive empirical functional --- adapts automatically to the local H\"{o}lder smoothness $\gamma > 1/2$ of $F_0$ at the estimation point, without knowing $\gamma$. We also prove a matching minimax lower bound, showing that no estimator can do better. The IIE is therefore the unique efficient nonparametric estimator in this problem. When $F_0$ is locally flat, the IIE even attains the faster $\sqrt{n}$ rate automatically.

Chapter~3 constructs the Bayesian counterpart via a Dirichlet Process prior on the observable distribution and isotonisation of the posterior. We prove semiparametric Bernstein-von Mises theorems for the resulting IIP in both Wicksell's problem and the deconvolution problem on the half-line. The credible sets derived from the IIP have asymptotically correct frequentist coverage at the optimal rate $\sqrt{n}/\sqrt{\log n}$, adapting to $\gamma$ without estimating it. 

Chapter~4 extends the IIP methodology to general half-line deconvolution, where the forward operator is a convolution with a known monotone kernel. We derive both a Bernstein-von Mises theorem for the IIP and an asymptotic distribution result for the frequentist isotonized estimator, providing the first nuisance-parameter-free UQ results for this class of problems.

Chapter~5 turns to causal inference. We introduce a causal model for proximal inference based on mixtures of linear models for the regression of the outcome on an unmeasured confounder. Under this model the bridge function, which links observed proxies to the causal effect, exists and has a closed-form expression, making the average treatment effect identifiable in a richer class of settings than was previously known. We propose a Bayesian nonparametric estimator based on a probit stick-breaking process and prove it is consistent.


\chapter{Samenvatting}

In veel wetenschappelijke problemen is de grootheid die men wil schatten niet direct waarneembaar. In plaats daarvan zijn de gegevens een ruisig, indirect beeld daarvan, geproduceerd door een bekende voorwaartse transformatie. Het terugwinnen van de onbekende grootheid uit de waargenomen gegevens vereist het omdraaien van die transformatie: een statistisch invers probleem. Een kanoniek voorbeeld is Wicksells stereologische probleem: bolvormige deeltjes van onbekende grootte zijn ingebed in een ondoorzichtig materiaal. De enige waarneembare gegevens zijn de stralen van de cirkelvormige doorsneden die door een snijvlak worden achtergelaten, welke een bekende meetkundige transformatie zijn van de werkelijke bolstralen. Dezelfde wiskundige structuur komt voor in de astronomie, deconvolutieproblemen in de epidemiologie, en, in een andere gedaante, in causale gevolgtrekking met ongemeten verstorende variabelen.

Twee moeilijkheden maken deze problemen lastig. De eerste is \textit{slecht-gesteldheid}: de voorwaartse operator kan de gegevens afvlakken, waardoor het inverteren ervan schattingsfouten versterkt; naïeve onbeperkte schatters zijn grillig en hebben vaak oneindige variantie. Het gebruikmaken van \textit{vormrestricties} zoals monotoniciteit, in sommige van de hier behandelde problemen, is de natuurlijke regularisatie. De tweede moeilijkheid is \textit{onzekerheidskwantificering} (UQ). Zelfs na het construeren van een goed-gedragende schatter is het bouwen van geldige betrouwbaarheidsgebieden niet triviaal: de limietverdelingen hangen af van onbekende lokale eigenschappen van de verdeling $F_0$ (de dichtheid van de waarnemingen $g_0(x)$ en de lokale Hölder-gladheid $\gamma$ van $F_0$), die op zichzelf moeilijk te schatten zijn. Dit maakt naïeve insteek-betrouwbaarheidsintervallen onbetrouwbaar.

De voornaamste aanpak die in dit proefschrift wordt ontwikkeld, is het construeren van Bayesiaanse niet-parametrische posteriorverdelingen en vervolgens het isotoniseren van de posteriortrekkingen --- door elke trekking te projecteren op de verzameling van monotone functies via de kleinste concave majorant. Het resultaat, de \textit{Geïsotoniseerde Inverse Posterior} (IIP), is computationeel goedkoop en, zoals wij aantonen, automatisch geldig voor onzekerheidskwantificering zonder dat het schatten van enige parameter vereist is.

Hoofdstuk~2 bestudeert Wicksells probleem vanuit een frequentistisch perspectief. Wij bewijzen dat de \textit{isotonische inverse schatter} (IIE) --- de rechter afgeleide van de kleinste concave majorant van een naïeve empirische schatter --- zich automatisch aanpast aan de lokale Hölder-gladheid $\gamma > 1/2$ van $F_0$ in het schattingspunt, zonder $\gamma$ te kennen. Wij bewijzen ook een overeenkomende minimax-ondergrens, waaruit blijkt dat geen enkele schatter het beter kan doen. De IIE is daarmee de uniek efficiënte niet-parametrische schatter in dit probleem. Wanneer $F_0$ lokaal vlak is, bereikt de IIE zelfs automatisch de snellere $\sqrt{n}$-snelheid.

Hoofdstuk~3 construeert het Bayesiaanse analogon via een Dirichlet-procesprior op de waarneembare verdeling en isotonisering van de posterior. Wij bewijzen semiparametrische Bernstein-von Mises-stellingen voor de resulterende IIP in zowel Wicksells probleem als het deconvolutieprobleem op de halve rechte lijn. De geloofwaardigheidsverzamelingen afgeleid uit de IIP hebben asymptotisch correcte frequentistische dekking op de optimale snelheid $\sqrt{n}/\sqrt{\log n}$, met aanpassing aan $\gamma$ zonder deze te schatten.

Hoofdstuk~4 breidt de IIP-methodologie uit naar algemene deconvolutie op de halve rechte lijn, waarbij de voorwaartse operator een convolutie is met een bekende monotone kern. Wij leiden zowel een Bernstein-von Mises-stelling voor de IIP als een asymptotisch verdelingsresultaat voor de frequentistische geïsotoniseerde schatter af, en leveren daarmee de eerste resultaten voor onzekerheidskwantificering zonder parameter voor deze klasse van problemen.

Hoofdstuk~5 richt zich op causale gevolgtrekking. Wij introduceren een causaal model voor proximale inferentie gebaseerd op mengsels van lineaire modellen voor de regressie van de uitkomst op een ongemeten verstorende variabele. Onder dit model bestaat de brugfunctie, die waargenomen proxies koppelt aan het causale effect, en heeft zij een expliciete uitdrukking, waardoor het gemiddelde behandelingseffect identificeerbaar is in een rijkere klasse van situaties dan voorheen bekend was. Wij stellen een Bayesiaanse niet-parametrische schatter voor gebaseerd op een probit-stokbreekproces en bewijzen dat deze consistent is.